\section{Three constituents in a five-dimensional module}
\label{fl221:sec:geometry}

A filtration with successive dimensions $(2,2,1)$ can be assembled
in two orders. One first joins the two planes and then adds the line.
The other first joins the second plane to the line. The second order
requires a coefficient map involving a three-dimensional constituent.
The comparison concerns these maps before taking cohomology.

Extend the notation to $1\le n\le5$:
\[
 V_n=\Mat_n(\C)^3,\qquad f_n(A,B,C)=\Tr(A[B,C]),\qquad
 Z_n=\{[A,B]=[B,C]=[C,A]=0\}.
\]
The trace pairing gives the displayed critical equations. Write
$\frX_n=[V_n/\GL_n]$ and $\frM_n=[Z_n/\GL_n]$.
All coefficients are rational. Complexes on a quotient retain the
whole action diagram and its composition identifications. Ordinary
atlas pullback carries no extra shift. The coefficient is
\[
 \cP_n=(\Phi_{f_n}\Q_{\frX_n}[2n^2])|_{\frM_n},
 \qquad \Phi_f=\Cone(i^*\longrightarrow\psi_f)[-1].
\]
Shifts and the tensor differential follow the preceding chapter.
In particular $\cP_1=\Q[2]$. The product order throughout is
submodule followed by quotient.

Let $F=\operatorname{Fl}(2,4;5)$ be the variety of flags
$H\subset K\subset\C^5$ with dimensions two and four. Let
$E_{221}\subset V_5\times F$ classify triples preserving this flag.
The three successive quotients have dimensions $(2,2,1)$.
In a basis adapted to the flag, all matrices are upper block triangular.
Multiplying these matrices gives
\[
 f_5|_{E_{221}}=f_2(A_H,B_H,C_H)
                  +f_2(A_{K/H},B_{K/H},C_{K/H}).
\]
Both summands are nonzero functions. For example, the triple
$(\operatorname{diag}(1,-1),e_{12},e_{21})$ has potential two.
Two such diagonal blocks give value four. Reversing the first
matrix of the second block gives values two and minus two.
The zero fibre of the sum therefore contains points outside the
product of zero fibres.

Forgetting $H$ gives a proper refinement
$r_L:E_{221}\to E_{4,1}$. Forgetting $K$ gives
$r_R:E_{221}\to E_{2,3}$. Their potentials remain
\[
 f_5|_{E_{4,1}}=f_4|_K+0,\qquad
 f_5|_{E_{2,3}}=f_2|_H+f_3|_{\C^5/H}.
\]
On the right refinement the second term pulls back to
$f_2|_{K/H}+0$. It remains the genuine function $f_3$ on the
unrefined target.

The dimension of $F$ is $4\cdot1+2\cdot2=8$.
A preserving matrix has $4+4+1+4+2+2=17$ entries. Hence
\[
 \dim E_{221}=3\cdot17+8=59,\qquad \dim V_5=75.
\]
The lower blocks $(21),(31),(32)$ have sizes four, two and two.
There are twenty-four forbidden entries for the three matrices.
The normal bundle of $E_{221}\subset V_5\times F$ has complex
rank twenty-four. Its positive complex Thom class has degree
forty-eight. Integration along $F$ has degree minus sixteen.
Their composition is
\[
 \gamma_{221}:Rq_{221*}\Q_{E_{221}}\longrightarrow\Q_{V_5}[32].
\]
Here $q_{221}$ forgets the flag. This is a proper map because $F$
is projective. The complex orientation fixes the sign.

\section{The coefficient maps needed by the two orders}
\label{fl221:sec:coefficient}

A morphism suffices to pass coefficients through an incidence.
Its construction can precede a proof of its invertibility.
For analytic functions $f,g$, set
$P_f=\{\Re f\le0\}$ and $P_g=\{\Re g\le0\}$.
The relative-cochain product followed by the support inclusion gives
\[
 \TS_{f,g}:\Phi_fA\boxtimes\Phi_gB
 \longrightarrow k^*\Phi_{f+g}(A\boxtimes B),\qquad
 k:f^{-1}(0)\times g^{-1}(0)\hookrightarrow(f+g)^{-1}(0).
\]
The intermediate support complexes are real constructible complexes.
Their restrictions giving vanishing cycles are complex constructible.
The external-support construction and its proof are given in
Lemma~\ref{ts22:lem:support-product}. They apply to arbitrary
constructible coefficients. The half-plane characterization uses
cochains of a tube relative to its positive half-tube. Shrinking
these neighborhoods gives the restriction maps of nearby cycles.
Thus the support morphism uses the normalized $\Phi$, including
for $f_3$ and $f_4$, without evaluating their scalar Milnor fibres.

For clarity, its required naturality has a direct cochain meaning.
A map of coefficient complexes induces a map of the relative
complexes. A proper map identifies relative cochains near a compact
fibre with cochains on inverse images of shrinking neighborhoods.
A finite subcover proves the required cofinality. These identifications
commute with the product of relative cochains and with inclusions
of supports. Consequently, for proper $q:Y\to X$,
\[
 \Phi_gRq_*B\simeq Rq_{0*}\Phi_{gq}B.
\]
Here $q_0$ is the map between zero fibres.
The two triple products start with support
$P_f\times P_g\times P_h$ and end with support $P_{f+g+h}$.
Their intermediate inclusions factor the same inclusion.
The product-cell associator preserves their ordered cochains.
This proves equality of these morphisms. It asserts no invertibility
for a comparison involving $f_3$ or $f_4$.

For the pairs $(d,e)=(2,2),(2,1),(4,1),(2,3)$, let
$\cE_{d,e}=[E_{d,e}/\GL_{d+e}]$ be the ambient incidence.
Its block map $a^{\mathrm{amb}}_{d,e}$ is smooth. In block frames,
its upper entries form an affine space of dimension $3de$.
Changes of splitting form an additive group of dimension $de$.
The relative stack dimension is therefore $2de$. Its potential is
the pullback of $f_d+f_e$.

Put $n=d+e$ and retain the complete critical correspondence
\[
 \cF_{d,e}=\cE_{d,e}\times_{\frX_n}\frM_n,
 \qquad
 \frM_d\times\frM_e\xleftarrow{a_{d,e}}\cF_{d,e}
                       \xrightarrow{b_{d,e}}\frM_n.
\]
Commuting triples induce commuting triples on the submodule and
quotient, which defines $a_{d,e}$. The equations on the upper
entries remain in $\cF_{d,e}$. Smoothness is used for the ambient
map only. For example, take $A=\operatorname{diag}(0,0,1,1,2)$,
$B=e_{13}$ and $C=0$, with the standard coordinate flag.
All diagonal blocks commute, whereas $[A,B]=-e_{13}$.
This point is excluded by the full critical fibre product.

The complex Thom map of the three forbidden lower blocks, followed
by Grassmannian integration, is
\[
 \gamma_{d,e}:Rq_{d,e*}\Q_{E_{d,e}}\longrightarrow\Q_{V_n}[4de].
\]
Its normal degree is $6de$ and its fibre degree is $-2de$.
Apply $\TS_{f_d,f_e}$, smooth ambient pullback, and then restrict to
$\cF_{d,e}$. This gives the coefficient map into
\[
 \left.\Phi_{f_n|\cE_{d,e}}
       \Q_{\cE_{d,e}}[2d^2+2e^2]\right|_{\cF_{d,e}}.
\]
The zero-fibre restriction in $k$ is legitimate: both diagonal blocks
are critical. Proper exchange and
$\Phi_{f_n}(\gamma_{d,e}[2d^2+2e^2])$ then define
\[
 \mu_{d,e}:Rb_{d,e*}a_{d,e}^*(\cP_d\boxtimes\cP_e)
                             \longrightarrow\cP_n.
\]
The shift identity $2d^2+2e^2+4de=2n^2$ gives degree zero.

For $(2,2)$ this is the map of Section~\ref{ts22:sec:product}.
For $(2,1)$ and $(4,1)$ one potential is zero and the external
comparison is the product identification. For $(2,3)$ the displayed
construction uses $\TS_{f_2,f_3}$ as a morphism. No inverse to it
enters the definition of $\mu_{2,3}$.

\section{One flag and two factorizations}
\label{fl221:sec:flag}

Write $\cE_{221}=[E_{221}/\GL_5]$ and
$\cF_{221}=\cE_{221}\times_{\frX_5}\frM_5$.
Its maps are $b_{221}:\cF_{221}\to\frM_5$ and
$a_{221}:\cF_{221}\to\frM_2\times\frM_2\times\frM_1$.
Set
\[
 \mathcal A=\cP_2\boxtimes\cP_2\boxtimes\cP_1,\qquad
 \mathcal D=Rb_{221*}a_{221}^*\mathcal A.
\]
The triple support product, equivalently
$\TS_{f_2,f_2}$ followed by the zero-factor comparison, gives
\[
 a_{221}^*\mathcal A\longrightarrow
 \left.\Phi_{f_5|\cE_{221}}\Q_{\cE_{221}}[18]
                                             \right|_{\cF_{221}}.
\]
The incoming shift is $8+8+2=18$. The known rank-two comparison
makes this arrow an isomorphism, although only the displayed map
will be used. Proper exchange and $\gamma_{221}[18]$ give
\[
                      M_{221}:\mathcal D\longrightarrow\cP_5.
\]
Its total shift is $18+32=50$.

The refinement squares on critical stacks are
\[
\begin{array}{ccc}
 \cF_{221}&\longrightarrow&\cF_{2,2}\times\frM_1\\
 \downarrow&&\downarrow b_{2,2}\times1\\
 \cF_{4,1}&\xrightarrow{a_{4,1}}&\frM_4\times\frM_1
\end{array}
 \qquad
\begin{array}{ccc}
 \cF_{221}&\longrightarrow&\frM_2\times\cF_{2,1}\\
 \downarrow&&\downarrow1\times b_{2,1}\\
 \cF_{2,3}&\xrightarrow{a_{2,3}}&\frM_2\times\frM_3.
\end{array}
\]
The left square adds an invariant plane inside $K$.
The right square adds an invariant plane inside $\C^5/H$.
These are exactly the data of the original flag, including its
frame changes and off-diagonal commutator equations. Thus both
squares are Cartesian as quotient stacks.

Proper base change and the external cochain product give
identifications with the two iterated sources:
\[
\begin{split}
 C_L:\mathcal D&\xrightarrow{\sim}
 Rb_{4,1*}a_{4,1}^*
  \bigl((Rb_{2,2*}a_{2,2}^*(\cP_2\boxtimes\cP_2))
                                      \boxtimes\cP_1\bigr),\\
 C_R:\mathcal D&\xrightarrow{\sim}
 Rb_{2,3*}a_{2,3}^*
  \bigl(\cP_2\boxtimes
       (Rb_{2,1*}a_{2,1}^*(\cP_2\boxtimes\cP_1))\bigr).
\end{split}
\]
These identifications use sheaf-theoretic proper base change. They
require no Tor-independence assertion about scheme intersections.
Define the two parenthesized maps by
\[
\begin{split}
 M_L&=\mu_{4,1}\circ
 Rb_{4,1*}a_{4,1}^*(\mu_{2,2}\boxtimes1)\circ C_L,\\
 M_R&=\mu_{2,3}\circ
 Rb_{2,3*}a_{2,3}^*(1\boxtimes\mu_{2,1})\circ C_R.
\end{split}
\]
Their shifts are respectively
\[
        18\longrightarrow34\longrightarrow50,
        \qquad18\longrightarrow26\longrightarrow50.
\]

\begin{theorem}[Associativity for two planes and a line]
\label{fl221:thm:assoc}
The maps just constructed satisfy
\[
                         M_L=M_{221}=M_R
\]
in the rational constructible derived category on $\frM_5$.
They retain the conjugation action, potential monodromy and
specialization after pullback to any analytic family of commuting
triples. The assertion uses no isomorphism for
$\TS_{f_2,f_3}$ and no scalar calculation for $f_3$ or $f_4$.
\end{theorem}

\begin{proof}
There are two parts: the normal orientations and the coefficient maps.
For the left refinement, the inner $(2,2)$ incidence contributes
twelve complex normal directions. The outer $(4,1)$ incidence
contributes twelve more. For the right refinement, the inner $(2,1)$
incidence contributes six and the outer $(2,3)$ contributes eighteen.
Both lists partition the same twenty-four forbidden matrix entries.
A filtration of a complex normal bundle has the product Thom class:
in local coordinates it is the product of the oriented normal disks.
The transition matrices are complex linear and preserve this orientation.
Each real Thom degree is even, so regrouping contributes sign $+1$.

The flag variety is a $\operatorname{Gr}(2,4)$-bundle over
$\operatorname{Gr}(4,5)$ on the left. It is a
$\operatorname{Gr}(2,3)$-bundle over $\operatorname{Gr}(2,5)$ on
the right. Integration in either order evaluates on its complex
fundamental cycle. The fibre dimensions are $4+4=2+6=8$.
Thus, on ambient representation spaces,
\[
 \gamma_{4,1}\,Rq_{4,1*}\gamma_{\mathrm{ref},L}
       =\gamma_{221}
       =\gamma_{2,3}\,Rq_{2,3*}\gamma_{\mathrm{ref},R}.
\]
The Gysin degrees are $16+16=8+24=32$.
Smooth ambient block charts identify the refinement maps with
$\gamma_{2,2}$ and $\gamma_{2,1}$, respectively.

The right coefficient comparison needs particular attention.
Put $B=Rq_{2,1*}\Q_{E_{2,1}}$ on $V_3$.
Its potential is $f_3$, and proper exchange identifies
$\Phi_{f_3}B$ with $Rq_{2,1,0*}\Phi_{f_2+0}\Q_{E_{2,1}}$.
The function $f_2+0$ is nonzero. Apply support naturality to
\[
                 \gamma_{2,1}:B\longrightarrow\Q_{V_3}[8].
\]
On $V(f_2)\times V(f_3)$ the square is
\[
\begin{array}{ccc}
 \Phi_{f_2}\Q\boxtimes\Phi_{f_3}B
 &\xrightarrow{\TS_{f_2,f_3}}&
 k^*\Phi_{f_2+f_3}(\Q\boxtimes B)\\
 \downarrow1\boxtimes\Phi\gamma_{2,1}&&
 \downarrow k^*\Phi(1\boxtimes\gamma_{2,1})\\
 \Phi_{f_2}\Q\boxtimes\Phi_{f_3}\Q[8]
 &\xrightarrow{\TS_{f_2,f_3}}&
 k^*\Phi_{f_2+f_3}(\Q\boxtimes\Q)[8].
\end{array}
\]
The top row is defined with coefficient $B$ and the bottom with
constant coefficients. Proper exchange identifies the top row with
the pushforward of $\TS_{f_2,f_2+0}$ on $V_2\times E_{2,1}$.
Both routes around the square are the product of relative cochains,
the same coefficient morphism, and the same inclusion of supports.
This proves its commutativity as a derived morphism. Neither
$\Phi_{f_3}B$ nor the target potential has been replaced by zero.

On the left apply the same naturality to
$Rq_{2,2*}\Q_{E_{2,2}}\to\Q_{V_4}[16]$, with the other factor
of rank one. The source potential is $f_2+f_2$, and the target
potential is $f_4$. Triple support associativity identifies the
common incoming map with the one defining $M_{221}$.
Commuting these squares with the ambient Gysin identity, then
applying proper exchange and restriction to $\frM_5$, gives
$M_L=M_{221}=M_R$.

All constructions commute with changes of frame on the action
diagrams. The relative product and its product-cell homotopies
commute with successive pullbacks. Thom classes and flag cycles
have their complex orientations on every level. The equalities
therefore descend with their composition identifications and
retain the stabilizers.

The associator retains the order of the three inputs. Every shift
in the construction is even, so it introduces no further sign.
Rotating the common support half-plane commutes with all its
inclusions and with the coefficient maps. This proves compatibility
with potential monodromy. Pull back the equality to an analytic
family and apply the natural map from the central restriction to
nearby cycles. This gives specialization compatibility for the
pulled-back ambient coefficient complexes. It does not replace them
by vanishing cycles of the identically zero function on the family.
\end{proof}

The ordinary cohomological operation $m_{d,e}$ is external product,
pullback by $a_{d,e}$, and application of $\mu_{d,e}$.
The theorem gives, for homogeneous critical cohomology classes,
\[
 m_{4,1}(m_{2,2}(x,y),z)=m_{2,3}(x,m_{2,1}(y,z)),
 \quad x,y\in H^*(\frM_2,\cP_2),\quad
 z\in H^*(\frM_1,\cP_1).
\]
No permutation occurs even when both $x$ and $y$ are odd.

\section{Regular matrices and a fivefold collision}
\label{fl221:sec:cyclic}

A cyclic matrix admits a vector whose first $n$ iterates form a
basis. Its commuting matrices are polynomials in it: a commuting
map is determined by the image of that vector. Its centralizer
therefore has dimension $n$ and is commutative.

\begin{lemma}[The coefficient over a cyclic matrix]
\label{fl221:lem:cyclic}
For $1\le n\le5$, on the critical locus where $A$ is cyclic,
the coefficient $\cP_n$ is the constant complex $\Q[2n]$.
The trivialization $u_n$ uses the complex orientation of the
maximal isotropic $B$ normal bundle and is compatible with conjugation.
The generators $e_n=\epsilon_nu_n$, with
$\epsilon_n=(-1)^{\binom n2}$, define a second trivialization.
\end{lemma}

\begin{proof}
The exact expression for the function is
\[
                         f_n=\Tr(B[C,A]).
\]
The linear map $C\mapsto[C,A]$ has rank $r=n^2-n$.
Its kernel is the centralizer of $A$. The kernel of the transpose
under the trace pairing is the same centralizer for $B$.
On the open cyclic locus these kernels form vector bundles.
Choose local complements. The function becomes exactly
$\sum_{i=1}^r u_i v_i$, with all centralizer coordinates free.
Both $B$ and $C$ are critical precisely when they belong to their
centralizers. They then commute with each other. Thus these free
coordinates describe the full critical locus.

The nondegenerate quadratic has $2r$ complex variables. Its Milnor
fibre retracts to $S^{2r-1}$ by the real-imaginary-coordinate
retraction in the preceding chapter. Normalized vanishing cycles
have degree $2r$, so the shift $[2n^2]$ gives $\Q[2n]$.
The maximal isotropic $B$ directions have their complex orientation.
Use their positive relative Thom class and its cone representative
as in Section~\ref{ts22:sec:normalization}.
A change of complementary frames preserves this orientation and
the resulting generator. These local generators therefore glue
and are conjugation invariant. For $n=1$ the function is zero,
and the stated coefficient is directly $\Q[2]$.
\end{proof}


The comparison extends across the entire cyclic locus, including
nondiagonalizable matrices. The quotient of the $B$ bundle by the
centralizer bundle is a complex vector bundle of rank $n^2-n$.
Its dual is paired by $\Tr(B[C,A])$ with the corresponding $C$ quotient.
Local complements form affine spaces, and Hermitian metrics form
contractible spaces. They give the positive disks used to define $u_n$.
Complex frame changes preserve their orientations and their relative
Thom classes. Thus the disk generators agree on overlaps and on
all levels of the conjugation action diagram.
Multiplication by the constant $\epsilon_n$ gives $e_n$ on that
same whole diagram. It requires no eigenline splitting at a collision.

Both trivializations respect potential monodromy. For a quadratic
in $2r$ variables, continuation once around zero induces the antipodal
map on $S^{2r-1}$, of degree $(-1)^{2r}=1$.
The zero-potential case $n=1$ has identity monodromy directly.
More generally the scalar automorphism $\epsilon_n1_{\cP_n}$
commutes with all continuation maps of the entire coefficient complex.

For any analytic family pulled back from the cyclic locus, the
central-restriction-to-nearby-cycles map commutes with these constant
scalar changes. On an ordered flag family with cyclic constituents,
the incoming scalar is the product of the constituent scalars.
The outgoing scalar is that of their total rank. Proper direct image
and the coefficient map retain this comparison. Consequently the
$u$ and $e$ descriptions give conjugate diagrams of actual sheaf maps
at both regular and singular fibres. Vanishing cycles remain the
pullback of the ambient coefficient complex.

When $A$ has five distinct eigenvalues, there are ten pairs of
eigenlines. Each pair contributes four quadratic normal variables.
Thus the full coefficient is $\Q[10]$. The input has the same
shift $[4+4+2]=[10]$. For each flag the complementary normal
quadratic has eight cross-eigenline pairs. Each lower-normal
projection has degree $-1$ in the $B$ polarization, so their product
has degree $(-1)^8=1$. Proposition~\ref{ts22:prop:normalization}
therefore gives coefficient $+1$ for each flag sheet in the $u$
trivializations of Lemma~\ref{fl221:lem:cyclic}.
The $e$ trivializations give the same coefficient because the incoming
change is $\epsilon_2^2\epsilon_1=1$ and the outgoing change is
$\epsilon_5=1$.

There are $\binom54\binom42=30$ flags in the left description
and $\binom52\binom32=30$ in the right. Both derived maps restrict
to the sum of these same thirty sheet coordinates. Conjugation invariance of
$u_n$ and $e_n$ makes this description independent of the eigenbasis.
Thus the trace normalization is the same in both specified conventions.

A collision retaining cyclicity gives more information than a sheet
count. Take a disk $D$ with coordinate $s$ and the companion matrix
\[
 A_s=\begin{pmatrix}
 0&1&0&0&0\\0&0&1&0&0\\0&0&0&1&0\\0&0&0&0&1\\s&0&0&0&0
 \end{pmatrix},\qquad B_s=C_s=0.
\]
Its characteristic polynomial is $t^5-s$, and it is cyclic even
at $s=0$. The central matrix is one nilpotent Jordan block.
It is not a scalar matrix.

\begin{proposition}[The flag cover at the collision]
\label{fl221:prop:cover}
The invariant $(2,2,1)$ flags over $D$ form a finite flat cover
$\pi:Y\to D$ of degree thirty. Its reduced analytic curve has six
smooth branches through one central point. Each branch maps by
$s=e^5$. The monodromy on a nearby fibre consists of six cycles
of length five.
\end{proposition}

\begin{proof}
A cyclic module is $\C[t]/(t^5-s)$. An invariant flag is equivalent
to ordered monic factors of degrees two, two and one. This statement
also holds over a local coefficient algebra: the cyclic quotient
has its monic characteristic polynomial, and polynomial division
identifies the kernel with its generated ideal. Applying division
successively gives the flag and its free successive quotients.

Write the linear factor as $t-e$. Then $e^5=s$. The other factors
satisfy
\[
 (t^2+at+b)(t^2+ct+d)=t^4+et^3+e^2t^2+e^3t+e^4.
\]
Their first two coefficients give
$c=e-a$ and $d=e^2-b-ae+a^2$. The remaining equations are
\[
\begin{split}
 g_1&=a^3-a^2e-2ab+ae^2+be-e^3=0,\\
 g_2&=a^2b-abe-b^2+be^2-e^4=0.
\end{split}
\]
The algebra is free over $\C[e]$ with basis
$1,a,a^2,b,ab,b^2$. Here is a direct division proof.
Adjoin
\[
\begin{split}
 g_3&=ag_2-bg_1=ab^2-ae^4-b^2e+be^3,\\
 g_4&=ag_3-bg_2=b^3-a^2e^4+abe^3-b^2e^2+be^4.
\end{split}
\]
Order monomials in $a,b$ by $(\deg_a+2\deg_b,-\deg_b)$,
with $e$ a coefficient. Their leading monomials are
$a^3,a^2b,ab^2,b^3$, all with coefficient one.
The adjacent leading-term cancellations are
\[
 bg_1-ag_2=-g_3,\qquad bg_2-ag_3=-g_4,\qquad
 bg_3-ag_4=e^4g_1-e^3g_2+e^2g_3-eg_4.
\]
Every term on the right has smaller leading monomial than the
corresponding cancellation. Nonadjacent cancellations are telescoping
sums of these adjacent ones. Division therefore gives a unique
remainder in the six listed monomials. No division by $e$ occurs.
Adjoining $e^5=s$ makes the algebra free of rank thirty over $\C[s]$.
At $s=0$ its only geometric point has $e=a=b=0$.

For $e\ne0$, put $\zeta=\exp(2\pi i/5)$. The four remaining roots
are $e\zeta^i$, $1\le i\le4$. Choosing two gives
\[
 a=-e(\zeta^i+\zeta^j),\qquad b=e^2\zeta^{i+j},\qquad i<j.
\]
These are six distinct smooth branches parameterized by $e$.
They meet only at zero. The algebra is generically reduced, and
its freeness over $\C[s]$ excludes any nilpotent supported at zero.
Thus these are its reduced analytic branches as well.
Following a loop in $s$ sends $e$ to $\zeta e$ on each branch,
which proves the asserted six five-cycles.
\end{proof}

For this family the input and output $u$ and $e$ trivializations agree:
$e_2\boxtimes e_2\boxtimes e_1=u_2\boxtimes u_2\boxtimes u_1$
and $e_5=u_5$. This equality holds over the central point as well,
by the cyclic construction above. The following trace uses these
identical trivializations.

\begin{proposition}[Coefficient specialization at the collision]
\label{fl221:prop:gluing}
Pulling $M_L=M_{221}=M_R$ back to this disk gives one morphism
\[
                   R\pi_*\Q_Y[10]\longrightarrow\Q_D[10].
\]
On $D\setminus\{0\}$ it is the thirty-sheet trace. At zero it is
multiplication by thirty. Let $j:D\setminus\{0\}\hookrightarrow D$
and $i:\{0\}\hookrightarrow D$ be the inclusions.
If $L$ is the rank-twenty-nine kernel of the sheet trace, then
\[
 R\pi_*\Q_Y=\Q_D\oplus j_!L,
 \qquad
 \dim H^0(i^*Rj_*L)=\dim H^1(i^*Rj_*L)=5.
\]
In particular $j_!L$ and $Rj_*L$ are different complexes.
\end{proposition}

\begin{proof}
The restrictions to the submodule and successive quotients are
cyclic modules with the three monic characteristic factors.
Lemma~\ref{fl221:lem:cyclic} identifies their coefficient product
with $\Q_Y[10]$ and the target with $\Q_D[10]$.
Proper base change gives the displayed actual coefficient morphism.
Finite fibres have no higher rational cohomology, so $R\pi_*\Q_Y$
is a sheaf in degree zero. Its central stalk is $\Q$, and its
nearby stalk is $\Q^{30}$. Specialization is the diagonal map:
a constant near the common branch point restricts to the same
constant on all thirty nearby points.

The target specialization is the identity. Compatibility of the
constructed morphism with specialization and the generic sheet
trace forces its central value to be thirty. This determines the
whole morphism: the target constant sheaf injects into its restriction
and direct image from the punctured disk, so a morphism vanishing
there is zero.

The constant-section map $\Q_D\to\pi_*\Q_Y$, followed by the
trace divided by thirty, is the identity. Its complementary kernel
has zero central stalk and restricts to $L$. The adjunction map
$j_!L$ to that kernel is an isomorphism on every stalk.

Let $W$ be a fibre of $L$ and $T$ its positive-circle monodromy.
The punctured disk retracts onto a circle. Its cellular cochains
are $W\xrightarrow{T-1}W$, in degrees zero and one.
A fixed vector in the thirty-sheet space is constant on each of
the six five-cycles. The trace-zero condition imposes one independent
relation on those six constants. Thus the kernel of $T-1$ on $W$
has dimension five. Its cokernel has the same dimension by
rank-nullity. This proves both displayed nearby-direct-image groups.
The stalk of $j_!L$ at zero is zero in every degree, proving the
last assertion.
\end{proof}

The five-cycle action is monodromy in the base parameter $s$.
Potential monodromy is a different action. On the cyclic normal
quadratic it is the antipodal map of $S^{39}$ and has degree $+1$.
The support construction intertwines the latter action and also
commutes with continuation in $s$. The trace calculation retains
both actions and their specialization maps.

\section{Scalar coefficients and the remaining geometric comparison}
\label{fl221:sec:scalar}

At a scalar five-dimensional triple every flag is invariant.
The incoming complex is the family
$\cP_2(H)\otimes\cP_2(K/H)\otimes\cP_1(\C^5/K)$
on $\operatorname{Fl}(2,4;5)$, with all frame actions retained.
At an ordinary atlas point, each rank-two scalar coefficient has
degrees $-4,-2,1$. The rank-one coefficient contributes degree
$-2$. Thus its nine ordered products have
\[
\begin{array}{c|rrrrrr}
 \text{degree}&-10&-8&-6&-5&-3&0\\\hline
 \text{rank}&1&2&1&2&2&1.
\end{array}
\]
The top product is the tensor of two odd degree-one classes and
the degree-minus-two line class. Interchanging the two plane
coefficients gives sign $-1$. Associating them in two orders
introduces no interchange and gives the equality proved above.

This is a calculation of incoming coefficient stalks. It does not
trivialize the family on the flag variety or determine its proper
direct image. The scalar target $\cP_5$, the image of $M_{221}$
there, and its extension data require further calculations.
The cyclic Jordan collision has a nonzero nilpotent matrix and
belongs to a different stratum.

For point modules, eigenvalues specify positions in $\C^3$ and
nilpotent parts specify extensions at coincident positions.
The two plane constituents retain cubic critical fluctuations,
including their odd scalar coefficient. The flag construction
couples those coefficients to the five-dimensional module in two
orders with the same result. This proves the stated rational
critical Hall associativity comparison. Additional block sizes,
higher coherence, a coproduct, Tate refinements and a compact
Calabi--Yau field-theory identification remain separate constructions.
